\section{Proof of \Cref{thm:pureucb}}\label{h1a:proof-pureucb}

In this section, we give a formal proof of our regret analysis of \PureUCB.

Let $x^*=\argmax_{x\in \calX} \min_{y\in \calY} u(x,y)$ and $y^* = \argmin_{y\in \calY} u(x^*,y)$ with ties broken arbitrarily. We know that $\vStar=u(x^*,y^*)$.

In each round, we know that $r_t$ is a noisy observation of $u(x!t,y!t)$. \Cref{code:ucb-n,code:ucb-sum,code:ucb-confidence} of \Cref{alg:pure-ucb} uses these observations to build a mean estimator for $u(x,y)$ and a confidence bound. Like the usual analysis of the UCB algorithm in stochastic MAB, we start by showing that $U!t$ is indeed an upper confidence bound on $u$, by showing the following lemma. We show its proof for completeness and for the extra details needed to handle the potential stochastic nature of the adversary.
\begin{lemma}\label{lem:ucb-ucb}
    Fix $(x,y)\in \calX\times \calY$. With probability at least $1-\delta$, the confidence bound $\absm[\Big]{u(x,y) - \frac{R!{t}(x,y)}{N!{t}(x,y)}}\leq \sqrt{\frac{4\log (1/\delta)+2\log (1+N!{t}(x,y))}{N!{t}(x,y)}}$ holds simultaneously for all $t$ such that $N!t(x,y)>0$.
\end{lemma}
\begin{proof}
    Denote as a shorthand $N_t=N!t(x,y)$ and $I_t=N_t-N_{t-1}$. Let $\calF_{t-1}$ be the $\sigma$-algebra generated by everything before $r_t$ is drawn, i.e., $(x!1,y!1,\dots,x!t,y!t,r_1,\dots,r_{t-1})$. By the UCB algorithm, $I_t$ and thus $N_t$ is measurable on $\calF_{t-1}$.
    
    By the definition of the noise model, the stochastic process $M=\{R!t(x,y)-N!t(x,y) u(x,y)\}_t$ is a martingale. Its difference sequence $\{M_t-M_{t-1}\}$ is conditionally 1-subgaussian when $(x!t,y!t)=(x,y)$, and is equal to 0 otherwise. Let $Z_t(\lambda)=\exp(\lambda M_t - \frac12 \lambda^2 N_t)$. Conditional on $\calF_{t-1}$, if $I_t=0$, then $Z_t(\lambda)=Z_{t-1}(\lambda)$; otherwise $Z_t(\lambda)=Z_{t-1}(\lambda)\exp(\lambda \eps_t - \frac12 \lambda^2)$. Combining both cases leads to $\EE\sbr{Z_t(\lambda)\mid\calF_{t-1}}\leq Z_{t-1}(\lambda)$, and thus $\{Z_t(\lambda)\}_t$ is a supermartingale.

    By the method of mixtures, $\Zbar_t=\int_\dR Z_t(\lambda) d\Phi(\lambda)$ is also a supermartingale, where $\Phi$ is the CDF of $\calN(0, 1)$. This integration evaluates to
    \begin{equation*}
        \Zbar_t=\int_\dR \exp(\lambda M_t - \frac12 \lambda^2 N_t) \frac{1}{\sqrt{2\pi}}\exp(-\frac 12 \lambda^2) d\lambda = \frac{1}{\sqrt{1+N_t}} \exp\rbrm[\Big]{\frac{M_t^2}{2(1+N_t)}}.
    \end{equation*}

    Since $Z_0(\lambda)=1$, we have $\Zbar_0=1$, and we can use Ville's inequality to get
    \begin{equation*}
        \PP\sbrm[\big]{\exists t: \Zbar_t\geq \frac{1}{\delta}}\leq \delta.
    \end{equation*}
    Plugging in, taking logarithms on both sides, and taking the complement on the event yields that with probability at least $1-\delta$, for all $t$,
    \begin{equation*}
        -\frac12 \log(1+N_t) + \frac{M_t^2}{2(1+N_t)} \leq \log(1/\delta),
    \end{equation*}
    which, when $N_t>0$, simplifies to
    \begin{equation*}
        \frac{|M_t|}{N_t}\leq \sqrt{\frac{2(1+N_t)}{N_t^2}\rbrm[\big]{\log(1/\delta)+\frac12 \log(1+N_t)}}.
    \end{equation*}
    Loosely bounding $\frac{1+N_t}{N_t}$ by 2 and plugging in our shorthands, we see that
    \begin{equation*}
        \frac{\absm[\big]{R!t(x,y)-N!t(x,y) u(x,y)}}{N!t(x,y)} \leq \sqrt{\frac{2}{N!t(x,y)}\rbrm[\big]{2\log(1/\delta)+\log(1+N!t(x,y))}},
    \end{equation*}
    which is equivalent to the desired result, up to algebraic simplification.
\end{proof}

Back to \Cref{thm:pureucb}. We first assume the high-probability event of \Cref{lem:ucb-ucb} happens for all action pairs, which implies $u(x,y)\leq U!t(x,y)$ for all $x$, $y$, and $t$. Using the confidence bound on the pair $(x^*,y)$, we see that
\begin{align*}
    \vStar & = u(x^*,y^*) 
    = \min_{y\in \calY} u(x^*,y)
    \leq \min_{y\in \calY} U!{t}(x^*,y) \\
    & \leq \min_{y\in \calY} U!{t}(x!t,y) \tag{$x!t$ maximizes $x\mapsto \min_{y\in \calY}U!{t}(x,y)$}  \\
    & \leq U!{t}(x!t,y!t). \yestag\label{eq:played-ucb-geq-vstar}
\end{align*}

Consider each action pair $(x,y)$ such that $\Delta_{xy}>0$. At the last iteration this action pair is played, i.e., at the maximum $t$ such that $(x!t,y!t)=(x,y)$, we know that $U!{t}(x!t,y!t)\geq \vStar=u(x,y)+\Delta_{xy}$, so
\begin{equation*}
    \sqrt{\frac{4\log (1/\delta)+2\log (1+N!{t-1}(x,y))}{N!{t-1}(x,y)}}\geq \Delta_{xy}.
\end{equation*}
which simplifies to $N!{t-1}(x,y)\leq \frac{1}{\Delta_{xy}^2} \rbrm[\big]{4\log (1/\delta)+2\log (1+N!{t-1}(x,y))}$, and can be loosely upper-bounded as $\frac{6}{\Delta_{xy}^2} \log T$. Thus, $N!t(x,y)$ is at most this value plus one. Since this is the last round $(x,y)$ is played, the same bound applies to $N!T(x,y)$.

Therefore, our regret is bounded as
\begin{align*}
    \psmr_T 
    & = \EE\sbrm[\Big]{\sum_{t=1}^T \rbrm[\big]{\vStar-u(x!t,y!t)}}
    = \EE\sbrm[\Big]{\sum_{(x,y)\in \calX\times \calY} N!T(x,y) \rbrm[\big]{\vStar-u(x,y)}} \\
    & = \EE\sbrm[\Big]{\sum_{(x,y)\in \calX\times \calY} N!T(x,y) \Delta_{xy}}
    \leq \EE\sbrm[\Big]{\sum_{(x,y) : \Delta_{xy}>0} N!T(x,y) \Delta_{xy}} \yestag\label{eq:ucb-wc-step}\\
    & \leq \sum_{(x,y) : \Delta_{xy}>0} \rbrm[\Big]{1+\frac{6}{\Delta_{xy}^2} \log T} \Delta_{xy}.
\end{align*}
This directly implies our instance-dependent bound.

To obtain our worst-case bound, we introduce a constant $\Delta>0$. For action pairs with $0<\Delta_{xy}<\Delta$, starting from \eqref{eq:ucb-wc-step}, we bound the total number of plays as $\sum_{(x,y):0<\Delta_{xy}<\Delta} N!T(x,y)\leq T$; for each action pair with $\Delta_{xy}>\Delta$, the regret contributed by such a pair is bounded by $\rbrm[\big]{1+\frac{6}{\Delta_{xy}^2} \log T} \Delta_{xy}\leq 2+\frac{6 \log T}{\Delta_{xy}}\leq 2+ \frac{6 \log T}{\Delta}$. Together, we have
\begin{align*}
    \psmr_T
    & \leq \Delta T + \sum_{(x,y) : \Delta_{xy}>\Delta} \rbrm[\Big]{2+ \frac{6 \log T}{\Delta}} \\
    & \leq \Delta T + m_x m_y \rbrm[\Big]{2+ \frac{6 \log T}{\Delta}} \\
    & = 2 m_x m_y + 2\sqrt{6 m_x m_y \log T},
\end{align*}
where in the last step we substitute in with $\Delta=\sqrt{\frac{6 m_x m_y \log T}{T}}$

Note that both results are conditional on the high-probability event in \Cref{lem:ucb-ucb} holding for all action pairs $(x,y)$, which happens with probability at least $1-m_x m_y \delta$. If the confidence event does not happen, we trivially bound the PSMR by $T$, and when $\delta=\frac{1}{T}$, this low-probability event adds an extra term of $m_x m_y$ to both bounds of PSMR; in both bounds, the extra term is absorbed by the big-$\order$ notation. \qed
